
A replication of \textit{"A Cyber-Physical System for Freeway Ramp Meter Signal Control Using Deep Reinforcement Learning in a Connected Environment"} from Hou
et al. (2021), published in the \textit{2021 IEEE Intelligent Transportation Systems Conference (ITSC)},
Indianapolis, USA, September 19--21, 2021.

\section*{Replication Summary}

\textbf{Scope of Replication} --- 
The objective of this work is to reproduce the results of the reference study \cite{hou2021rampmeter} and verify the claim that the reinforcement learning algorithms Ape-X DQN, PPO and A3C 
show superior performance for freeway ramp meter signal control in the metrics traffic throughput, vehicle speed, number of stops per vehicle, fuel efficiency per vehicle and per-vehicle emission rates for CO$_2$ and NO$_x$ 
than the given baseline controllers no ramp meter, fixed-time controller and ALINEA in both mixed near-congested/congested traffic situations and extremely congested traffic situations. 
Furthermore, we aim to confirm that PPO shows better results in all of the aforementioned assessment criteria than the other reinforcement learning algorithms, while showing a convergence trend similar to Ape-X DQN and one that is approximately 15 times faster than A3C. 
To achieve this, we implemented a simulation environment in Python, together with the three baseline controllers (no ramp meter, fixed-time control and ALINEA) and the three RL-based controllers (Ape-X DQN, PPO and A3C), which we subsequently used to compare our obtained results with the original claims. 

\vspace{1em}

\textbf{Methodology} --- 
Since the authors of the reproduced study did not provide any source code, the modelling of the environment, the route generation, the network, 
the reward function, the state computation, the three baseline controllers, the three reinforcement learning algorithms and also the training and evaluation scripts were implemented from scratch
based on the paper description. For Ape-X DQN, PPO and A3C we made use of the reinforcement learning library RLlib \cite{liang2018rllib} to remain consistent with the original study and the microscopic traffic simulations were conducted using the SUMO software (Simulation of Urban Mobility) together with TraCI (Traffic Control Interface) \cite{krajzewicz2012recent}.


\vspace{1em}

\textbf{Results} --- 
Our replication findings do not match the original work from \cite{hou2021rampmeter}. With their described parameters, the baseline controllers (no ramp meter, fixed-time ramp meter and ALINEA)
already showed notable deviations in the evaluation metrics. Furthermore, the hypothesis that the reinforcement learning policies outperform the baseline controllers could not be reproduced. 
Moreover, we were not able to provide evidence that PPO surpasses Ape-X DQN and A3C in our measurements. However, we can confirm that PPO converges at a similar rate as Ape-X DQN and that A3C takes significantly longer.  


\vspace{1em}

\textbf{What was easy} --- 
The implementation of the three baseline-feedback controllers (no ramp meter, fixed-time controller and ALINEA) was straightforward, especially no ramp meter and fixed-time controller. 
The reason for this is that the paper provided both formulas and parameter values for ALINEA, while also specifying the cycle length for the fixed-time controller. 

\vspace{1em}

\textbf{What was difficult} --- 
Implementing the environment with all its details required significant effort. 
